\scanpage{99}

déjà prouvé pour $n-1$, et que $s_0\ldots s_n$ et $(a_{\unsure{-1|-9}},\ldots,a_{n-1})$ sont lin. indép. et $(s_0,\ldots s_n)$ engendrent $\unsure{E_n|\mathcal E_n}$. \deleted{\unsure{Or on a|On a donc} \unsure{par la}}

Dire $(s_0,\ldots s_n,s_{n+1})$ libre \unsure{ssi} $s_{n+1}$ libre mod $E_n$ \deleted{$s_{n+1}\equiv\lambda_n a_n\bmod E_n$, donc $s_{n+1}$ \unsure{est|et} $(s_0\ldots s_{n+1})$ libre \unsure{ssi}} \deleted{$s_{n+1}$} or par \unsure{1°)|(6)} on a

\[s_{n+1}\equiv\lambda_n a_n\quad(E_n)\]

donc la condition équivaut à $\alpha)$ $a_n$ libre mod $E_n$, i.e. (18), et $\beta)$ \deleted{\unsure{$\lambda_n$}} la condition (19). \deleted{Pour}

Donc on a l'équivalence dite, et de plus ces \deleted{$E_n=k(s_0,\ldots,s_n$\unsure{$\ E_n$}} conditions équiv. impliquent

\[E_{n+1}=E_n+ks_{n+1}=k(s_0,\ldots,s_n)+ks_{n+1}=k(s_0,\ldots s_{n+1}),\]

ce qui achève la dém. du corollaire.

\textcircled{\scriptsize $2^\circ$} On procède par récurrence sur $n$, supposons (\unsure{si|ici} $n\ne0$) que c'est prouvé déjà pour $n-1$. Donc on peut supposer que ttes les conditions écrites dans A, B sont satisfaites pour le système $(\deleted{\unsure{E}}s_0,\tau_0,\ldots\tau_{n-1})$, \inserted{et qu'on a \deleted{\unsure{donc} $E_i=k(s_0\ldots s_{i-1}$} de plus (14) pour $0\leq i\leq n$,} et prouver qu'alors le syst. de conditions

\[(\mathrm A'')\quad\left\{\begin{aligned}\tau_n&\text{ commute à }\tau_0\ldots\tau_{n-2}&&\alpha).\\\tau_n s_0&=s_0 &&\beta)\\s_{n+1}&\text{ libre mod }E_n\ \deleted{\unsure{\}}}&&\gamma)\end{aligned}\right.\]

équivaut au syst. de conditions

\[(\mathrm B'')\quad\left\{\begin{aligned}\left\{\begin{aligned}\langle a_i,a'_n\rangle&=0\\ {}a'_i,a_n\rangle&=0\end{aligned}\right\}\quad\text{pour }0\leq i\leq n-2&&\big\}\ \alpha)\\ \langle s_0,a'_n\rangle=0&&\beta)\\ a_n\text{ libre mod }E_n&&\gamma)\\ \lambda_n\in k^*&&\delta)\end{aligned}\right.\]

et que ces conditions impliquent

\[E_{n+1}=k(s_0,\ldots s_{n+1}).\tag*{(20)}\]

% archive sheet number: 98
